\begin{textsample}{POS Dim 1 – mg – Score 27.00 – mg/mg\_ef\_linguagens\_lp\_9.txt}  \label{ex:f1_pos_248}
5.5.1 Multiletramentos A necessidade de expandir a visão de letramento tradicional, \textbf{focado} predominantemente nas habilidades de leitura e escrita da linguagem verbal ( JEWITT, 2008, p. 244 ), para \textbf{uma} proposta mais abrangente sobre a linguagem em \textbf{uma} perspectiva enunciativo-discursiva, considerando os variados recursos semióticos mobilizados, é consensual entre muitos \textbf{pesquisadores}.
Tal perspectiva da linguagem está relacionada ao contexto \textbf{contemporâneo} \textbf{que} é marcado, principalmente, por grandes desenvolvimentos tecnológicos e pelo fenômeno da globalização, o qual tem como \textbf{uma} de suas consequências a aproximação entre a imensa diversidade cultural e social.
Nesse contexto, apesar da relevância do letramento voltado para o exame da linguagem verbal, surge a necessidade de pensar novas formas de letramento. \textbf{Uma} das discussões iniciais sobre a expansão do conceito de letramento proposta pelo Grupo de Nova Londres ( 1996 ) apontou o termo “ multiletramentos ” para definir \textbf{uma} nova \textbf{abordagem}, a qual oferece argumentos para repensar os letramentos e suas implicações para a participação social na vida pública, econômica e comunitária ( THE NEW LONDON GROUP, 1996 ).
Várias atividades das quais participamos \textbf{exigem} o conhecimento de \textbf{inúmeros} saberes contemplados nos multiletramentos, \textbf{baseados} na manipulação de \textbf{uma} multiplicidade de linguagens, culturas, práticas sociais e contextos, diferentemente de \textbf{uma} visão de letramento \textbf{embasada} na apreensão de regras e sua aplicação de maneira correta ( COPE ; KALANTZIS, 2008 ).
À medida \textbf{que} nos \textbf{educamos} em \textbf{direção} aos multiletramentos, as ações em busca de \textbf{uma} participação mais influente na vida \textbf{contemporânea} se tornam mais informadas segundo conhecimentos e processos especializados.
Tal especialização está relacionada a saber interagir em situações ( gêneros discursivos ) familiares e não familiares e ser capaz de procurar por pistas para \textbf{uma} participação mais apropriada nessas práticas ( COPE ; KALANTZIS, 2008 ), \textbf{uma} vez \textbf{que} os saberes dos multiletramentos podem ser definidos como “ as habilidades de interagir com a pluralidade ( LO BIANCO, 2000, p. 99 ), como leitores e produtores de texto. ” ( MOTTA-ROTH ; HENDGES, 2010, p. 45 ). \textbf{Uma} vez \textbf{que} o objetivo pedagógico dos multiletramentos é “ proporcionar aos alunos \textbf{uma} \textbf{percepção} sobre como padrões de significação são produtos de diferentes contextos ” ( COPE, KALANTZIS, 2008, p. 205 ), o reconhecimento e uso desses padrões \textbf{depende}, dentre outros fatores, da manipulação de diferentes modos semióticos.
O letramento multimodal é \textbf{uma} das propostas da \textbf{pedagogia} dos multiletramentos, e está relacionado à referida manipulação de diferentes modos semióticos. 5.5.2 Letramento Multimodal Conforme indicado no manifesto de 1996 do Grupo de Nova Londres, \textbf{uma} das ideias principais \textbf{que} informa a noção de multiletramentos é a crescente complexidade e \textbf{inter-relação} de diferentes modos semióticos.
Dentre as grandes áreas \textbf{que} contemplam essa complexidade e \textbf{inter-relação}, o letramento multimodal recebe atenção especial por incorporar e reunir os saberes necessários para lidar com esses diversos modos semióticos ( THE NEW LONDON GROUP, 1996, p. 17 ).
Considerando os dois argumentos \textbf{que} justificam a multiplicidade defendida pelos multiletramentos, a diversidade cultural e linguística de \textbf{um} \textbf{lado}, e a \textbf{influência} de novas tecnologias comunicativas de outro, “ em ambos os casos, tem -se em \textbf{mente} modos de representação mais amplos e mais dinâmicos do \textbf{que} exclusivamente a linguagem verbal ” ( MOTTA-ROTH ; HENDGES, 2010, p. 45 ).
Tais constatações apontam para a importância do letramento multimodal, como \textbf{um} conjunto de práticas \textbf{que} consideram essa ampliação de foco da linguagem verbal para outros modos semióticos a fim de dar conta dessa multiplicidade.
Os saberes \textbf{que} envolvem os multiletramentos contemplam habilidades para interagir tanto com a diversidade de culturas e línguas quanto com a diversidade de tecnologias comunicativas.
Manipular as diversas tecnologias comunicativas define a proposta \textbf{central} do letramento multimodal, ou seja, conhecer o papel dos recursos semióticos e o uso integrado dos mesmos na construção de sentido.
Segundo Rojo ( 2012 ), existe \textbf{uma} necessidade de se incluir no currículo escolar a ampla multiplicidade das novas culturas e textos \textbf{que} surgem no mundo \textbf{globalizado} com o subsídio das novas tecnologias, \textbf{visto} \textbf{que} já estão presentes na vida dos estudantes, mas ainda não estão presentes na escola.
O ensino pautado nos textos multimodais possibilita ao estudante compreender essa diversidade encontrada nos mais diferentes meios de comunicação presentes no cotidiano, sendo, diante disso, cada vez mais “ frequente a preocupação dos professores em inserir gêneros textuais diversos e recursos tecnológicos da sociedade \textbf{moderna} nas atividades realizadas em sala de aula ” ( DIONÍSIO, 2005, p. 140 ).
E essa preocupação é norteada pelo seguinte fato :"todo professor tem convicção de \textbf{que} imagens ajudam a aprendizagem, quer seja como recurso para prender a atenção dos alunos, quer seja como portador de informação complementar ao texto verbal ” ( idem, p. 141, \textbf{grifo} da autora ).
Mayer ( 2001 ) apud Dionísio ( 2005 ) \textbf{diz} \textbf{que} “ os alunos aprendem melhor através de palavras e imagens \textbf{que} de palavras apenas ” ( p. 141 ) constate -se então \textbf{que} quando procuramos \textbf{uma} metodologia \textbf{que} abrange palavras e imagens, aliamos a \textbf{ludicidade} à interação e com isso \textbf{despertamos} \textbf{um} maior interesse nos estudantes em aprender.

% matched lemmas: abordagem, basear, central, contemporâneo, depender, despertar, direção, dizer, educar, embasar, exigir, focar, globalizado, grifo, influência, inter-relação, inúmero, lado, ludicidade, mente, moderno, pedagogia, percepção, pesquisador, que, um, ver
\end{textsample}
